\clearpage\SourceLeaf{118}\ArithmeticHead{33}{R}{III b. LEMNISKATISCHE FUNKTIONEN 1796—1799.}
\HistoricalFoot{X 2 Abh. 2.}{5}
An Aufzeichnungen kommen hier vorerst in Betracht:
\begin{list}{}{\setlength{\leftmargin}{3em}\setlength{\labelwidth}{1.5em}\setlength{\labelsep}{.5em}\setlength{\topsep}{0pt}\setlength{\itemsep}{0pt}\setlength{\parsep}{0pt}}
\item[1.] die Leistenotizen Werke X, 1, S. 152—164,
\item[2.] die Werke III, S. 404—406, Artikel 1—4 abgedruckten Fragmente,
die \textsc{Gauss} auf lose Blätter (Fh) etwa zu Anfang des Jahres 1798
aufgezeichnet haben dürfte.
\end{list}
\indent\textsc{Gauss} setzt:
\[
\begin{gathered}
x=\operatorname{sin\,lemn}\left(\int\limits_0^x\frac{dx}{\sqrt{1-x^4}}\right),\quad
\tfrac12\varpi=\int\limits_0^1\frac{dx}{\sqrt{1-x^4}},\\[1ex]
x=\operatorname{cos\,lemn}\left(\tfrac12\varpi-\int\limits_0^x\frac{dx}{\sqrt{1-x^4}}\right)\hbox{\footnotemark}
\end{gathered}
\]
\footnotetext{In der Handschrift bezeichnet \textsc{Gauss} den Wert des Integrals von 0 bis 1 mit \(\frac12\pi\) und die Umkehrungsfunktion mit
\[
x=\operatorname{sn}\int\frac{dx}{\sqrt{(1-x^4)}}=\operatorname{cs}\left(\tfrac12\pi-\int\frac{dx}{\sqrt{(1-x^4)}}\right).
\]
\indent In bezug auf die Bezeichnungen ist noch folgendes zu bemerken. An manchen Stellen (namentlich
nach 1798 in der Scheda Aa) hat \textsc{Gauss} für den sinus und cosinus lemniscaticus eigene Zeichen; 1797 im
\textit{Leiste} schreibt er anfangs dafür einfach sin, cos oder auch nur \(s,c\); aber auch die im Texte angewandten
(von \textsc{Schering} benutzten) Abkürzungen finden sich im \textit{Leiste} (z. B. bei S. 92, 93). Die Bezeichnung \(\varpi\) für
die halbe Periode gebraucht \textsc{Gauss} von 1798 an ganz konsequent. Der benutzte Buchstabe ist eine ältere,
den Druckwerken des XVIII. Jhs. eigentümliche Form des \(\pi\), die sich z. B. auch in der \textit{Introductio} von
\textsc{Euler} findet; \textsc{Gauss} will also daran erinnern, dass \(\varpi\) für die Lemniskate die analoge Rolle spielt, wie \(\pi\)
für den Kreis. In den Aufzeichnungen im \textit{Leiste} (Werke X, 1, S. 152—159) wird der vierte Teil der lemniskatischen
Periode mit \(\Pi\) oder geradezu mit \(\pi\) und in der Tagebuchnotiz 63. mit \(\pi^l\) (d. h. also das \(\pi\) der
Lemniskate!) bezeichnet.}
\noindent und stellt die Additions- und (ganzzahligen) Multiplikationsformeln zusammen.
Von den letzteren scheint \textsc{Gauss} alsbald zur Teilungsgleichung übergegangen
zu sein, denn am 19. März 1797 trägt er bereits in sein \textit{Tagebuch} die Notiz
60. ein: »Weshalb man bei der Teilung der Lemniskate in \(n\) Teile auf eine
Gleichung vom Grade \(n^2\) geführt wird«. Aus dieser Fassung geht auch unzweifelhaft
hervor, dass \textsc{Gauss} zu jener Zeit (März 1797) schon die zweite
(imaginäre) Periode entdeckt hatte, da die doppelte Periodizität, und nur diese,
den Grund dafür erkennen lässt, »\textls[120]{weshalb}« die Teilungsgleichung vom
Grade \(n^2\) ist\footnote{In seiner Gedächtnisrede auf \textsc{Jacobi} (\textsc{Jacobis} Werke I, S. 9) bemerkt \textsc{Dirichlet}: »So hatte
man .\ .\ . gefunden, dass der Grad der mit Hilfe des \textsc{Euler}schen Satzes gebildeten Gleichung, von deren
Lösung die Teilung des elliptischen Integrals abhängt, nicht wie in der analogen Frage der Kreisteilung
der Anzahl der Teile, sondern dem Quadrate dieser Anzahl gleich ist. Die Bedeutung der reellen Wurzeln,
deren Anzahl mit jener übereinstimmt, war leicht ersichtlich, wogegen die zahlreichen imaginären ganz unerklärlich
erscheinen mussten. Aber dass hier ein Geheimnis verborgen liege, darüber hatte man vor \textsc{Abel}
und \textsc{Jacobi} kein Bewusstsein, und ihnen war es vorbehalten, sich zuerst über diese und ähnliche Erscheinungen
zu wundern.« Vergl. auch a. a. O., S. 11.}. Damit war ihm aber nicht nur der Zugang zu den tiefsten\par
\clearpage\SourceLeaf{119}\ArithmeticHead{34}{L}{L. SCHLESINGER, ÜBER GAUSS’ FUNKTIONENTHEORETISCHE ARBEITEN.}
\noindent Ergebnissen der Theorie der lemniskatischen Funktionen frei gemacht, sondern
er fand sich auch genötigt, seine Funktionen \textls[120]{für komplexe Werte der
Variabeln} zu untersuchen.

Damit hängt wohl die Notiz zusammen, die \textsc{Gauss} auf der letzten Seite
des \textit{Tagebuchs} eingetragen hat (Werke X, 1, S. 515 bei Nr. 60 des \textit{Tagebuchs}):
»\textls[120]{Imaginäre Grössen}: Es wird ein allgemeines Kriterium verlangt, nach
dem man komplexe Funktionen von mehreren Variabeln von nichtkomplexen
unterscheiden kann«, und bei der das mit Bleistift geschriebene Datum: »1797,
Apr. 15.«, obwohl etwas verwischt, noch deutlich erkennbar ist. — Dass \textsc{Gauss}
die Bezeichnung »komplexe Grössen« erst verhältnismässig spät, öffentlich nicht
vor 1831, ständig benutzt, könnte es ja zweifelhaft erscheinen lassen, ob das,
was er hier unter komplexen, bezw. nichtkomplexen Funktionen versteht, sich
mit dem deckt, was wir nach dem heutigen Sprachgebrauch uns darunter
vorstellen; dabei muss allerdings bemerkt werden, dass in einer Aufzeichnung
(Werke X, 1, S. 472, siehe auch ebenda, S. 443), die aus dem Jahre 1805
stammen dürfte, der Ausdruck »Zug komplexer Werte« benutzt wird. Aber
wie immer es auch sei, so werden wir in dem Grad der Teilungsgleichung
das entscheidende Moment dafür zu erblicken haben, das \textsc{Gauss} veranlasst
hat, die lemniskatischen Funktionen für Werte des Arguments von der Gestalt
\(u+iv\) untersuchen, und zwar gleichgültig, ob wir annehmen, dass er die doppelte
Periodizität sofort in der uns heute geläufigen Form von Aggregaten
der beiden Perioden \(2\varpi\) und \(2i\varpi\) mit reellen ganzzahligen Koeffizienten
\(a\,.\,2\varpi+b\,.\,2i\varpi\), oder vielleicht zunächst in der Form von komplexen ganzzahligen
Vielfachen der einen Periode, also \((a+ib)2\varpi\) aufgefasst hat. Die
im \textit{Leiste} aufgezeichnete Tabelle, Werke X, 1, S. 154, Art. 5, gibt beiden
Auffassungen Raum\footnote{Sie gibt in der Ebene der komplexen Variabeln das quadratische \textls[120]{Punktgitter}, dessen Bedeutung
für die Theorie der biquadratischen Reste \textsc{Bachmann}, Abh. 1 dieses Bandes, S. 27, erörtert.}; man könnte sogar die zweite Auffassung, die sich den
arithmetischen Gedankengängen von \textsc{Gauss} anpasst, für die wahrscheinlichere\par
\clearpage\SourceLeaf{120}\ArithmeticHead{35}{R}{III b. LEMNISKATISCHE FUNKTIONEN 1796—1799.}
\HistoricalFoot{}{\(5^*\)}
\noindent halten, schon weil sich ihr das Auftreten der komplexen Multiplikation in
ungezwungenster Weise an die Seite stellen lässt (vergl. auch weiter unten
Abschn. IV b)).

Mit den unmittelbar folgenden Tagebuchnotizen 61., 62., 63. vom März
1797 korrespondieren die Aufzeichnungen im \textit{Leiste}, die Werke X, 1, S. 153
—163 abgedruckt sind (vergl. die zugehörigen Bemerkungen, S. 164—166),
die wir also mit Sicherheit für die hier in Rede stehende Zeit in Anspruch
nehmen können\footnote{Die Notiz 61. besagt, dass die Reihen
\[
S_k=\sum\left(\frac{m^4-6m^2n^2+n^4}{(m^2+n^2)^4}\right)^k,\qquad(k=1,\ 2,\ 3\ldots)
\]
wo die \(m,n\) alle Paare ganzer positiver Zahlen, mit Ausnahme von \(m=n=0\), durchlaufen, von den
Potenzen der lemniskatischen Periode \(\frac12\varpi\) abhängen (vergl. für den Beweis die Bemerkung Werke X, 1,
S. 516). \textsc{Gauss} hat also hier im lemniskatischen Falle die Klasse von Reihen betrachtet, die in allgemeinster
Form von G. \textsc{Eisenstein} (\textsc{Crelles} Journal 35, 1847, S. 153, Mathem. Abhandl. 1847, S. 213)
untersucht worden sind.}. \textsc{Gauss} bestimmt auf Grund der aus dem Additionstheorem
folgenden Formeln die Nullstellen und Unendlichkeitsstellen der Funktionen
sin lemn und cos lemn und leitet daraus die Darstellung dieser Funktionen
in der Form von Quotienten doppelt unendlicher Produkte her; auch für die
Konvergenzuntersuchung dieser Produkte findet sich ein Ansatz\begingroup\renewcommand{\thefootnote}{50a)}\footnote{So stehen im \textit{Leiste} (Werke X, 1, S. 154, 155) die Formeln
\[
\left(1-\frac{x^4}{\Pi^4}\right)\left(1-\frac{x^4}{16\Pi^4}\right)\left(1-\frac{x^4}{81\Pi^4}\right)\ldots,\quad
\sum\frac1{\rho^4}=2\left(\frac11+\frac1{16}+\frac1{81}+\text{etc.}\right).
\]
Und in der Scheda Aa (1798) heisst es (Werke III, S. 415): »Lässt man die Werte von \(\sin\varphi\), die \(\operatorname{sin\,lemn}\varphi
=0\) machen, nach bekannten Regeln ein unendliches Produkt erzeugen, ebenso die Werte von \(\sin\varphi\), die
\(\operatorname{sin\,lemn}\varphi=\infty\) machen, setzt das erstere gleich \(P\varphi\), das zweite gleich \(Q\varphi\), so wird es gestattet sein zu
setzen:
\[
\operatorname{sin\,lemn}\varphi=\frac{P\varphi}{Q\varphi},
\]
was wir strenge beweisen können.«}\endgroup\addtocounter{footnote}{-1}.
Die gedachten Funktionen erscheinen also als Quotienten von ganzen transzendenten
Funktionen, die \textsc{Gauss} anfangs mit \(M,N,\mu,\nu\) später mit \(P,Q,p,q\) bezeichnet.
Es werden nun diese ganzen Funktionen in gewöhnliche Potenzreihen der
Variabeln \(x\) entwickelt, und \textsc{Gauss} hebt (Werke III, S. 406) hervor, dass diese
Potenzreihen »schneller konvergieren als jede gegebene Konvergenz«, während
für die formal schon 1796 aufgestellte Potenzreihenentwicklung des sin lemn
jetzt festgestellt wird, dass sie divergiert, wenn die Variable grösser ist als\par
\clearpage\SourceLeaf{121}\ArithmeticHead{36}{L}{L. SCHLESINGER, ÜBER GAUSS’ FUNKTIONENTHEORETISCHE ARBEITEN.}
\noindent\(\dfrac{\varpi}{\sqrt2}\)\footnote{»Oder«, sagt \textsc{Gauss} gleich darauf schärfer, »wenn \(\varphi^4>\frac{\varpi^4}4\) gesetzt wird«; \(\varphi\) ist a. a. O. die unabhängige
Variable.}
und die für cos lemn, wenn die Variable grösser ist als \(\varpi\). Wie \textsc{Königsberger}\footnote{L. \textsc{Königsberger}, \textit{Zur Geschichte der Theorie der elliptischen Transzendenten}, 1879, S. 95.}
bemerkt hat, sind die Zähler und Nenner \(M,N\) genau die für den
lemniskatischen Fall spezialisierten Al-Funktionen von \textsc{Weierstrass}.

Bemerkenswert ist, dass \textsc{Gauss} im \textit{Leiste} (Werke X, 1, S. 153) neben
dem lemniskatischen Integral auch das Integral
\[
\int\frac{dx}{\sqrt{1-x^3}}
\]
betrachtet, offenbar weil es ihm lehrreich erschien, die bei den verschiedenen
Integralen auftretenden Realitätsverhältnisse miteinander zu vergleichen. Es
findet sich ebenda auch eine Darstellung der Umkehrfunktion dieses Integrals
als Quotient von doppelt unendlichen Produkten mit einem Ansatz zur Untersuchung
ihrer Konvergenz.

\textsc{Gauss} stellt nun (\textit{Leiste} S. 66—71, Werke X, 1, S. 156—158, Artikel 7, 8)
Formeln auf für die Multiplikation der Variabeln in den ganzen transzendenten
Funktionen \(M,N\) (den spätern \(P,Q\)) mit 2, 4 und mit der komplexen ganzen
Zahl \((1+\sqrt{-1})\), die fast vollständig mit denen übereinstimmen, die sich auf
S. 63 des Handbuchs 16, Bb, (begonnen 1801) finden und (in abgeänderter Reihenfolge)
Werke III, S. 410, 411 abgedruckt sind. Es sind dies zum Teil die
Ausführungen der in der Tagebuchnotiz 63. (vom 29. März 1797, Werke X, 1,
S. 517) unter 1) bis 3) angegebenen Resultate. In bezug auf die unter 4) und
5) ausgesprochenen Resultate, wonach für die Nennerfunktion \(N\) des sinus
lemniscaticus die Gleichung \(\log N(\varpi)=1{,}570796=\frac12\pi\) durch numerische Induktion
gefunden wurde, ist die entsprechende Rechnung Werke X, 1, S. 158,
Art. 9 zu vergleichen. \textsc{Gauss} sagt von dieser Gleichung in der Tagebuchnotiz,
sie sei »höchst bemerkenswert und ihr Beweis verspräche eine gewichtigste
Förderung der Analysis«. Wie wir sehen werden, gelingt ihm der Beweis
erst im Juli 1798 als Folge der Darstellung vom Zähler und Nenner des
sin lemn in der Form von einfach unendlichen Produkten.

Bei S. 78 des \textit{Leiste} (Werke X, 1, S. 167, Art. 2) stellt \textsc{Gauss} Differentialrelationen
zwischen den Funktionen \(P,Q\) (den frühern \(M,N\)) auf;\par
\clearpage\SourceLeaf{122}\ArithmeticHead{37}{R}{III b. LEMNISKATISCHE FUNKTIONEN 1796—1799.}
\noindent mit derselben Bezeichnung \(P,Q\) ist bei Seite 26 des \textit{Leiste} (Werke X, 1,
S. 167, Art. 1) ein Ausdruck für \(P(nx)\) für positiv ganzzahliges \(n\) durch
\(P,Q\) notiert. Wir haben hier einen interessanten Beleg dafür, dass \textsc{Gauss}
manchmal auf früher leer gelassenen Stellen seiner Notizhefte spätere Eintragungen
gemacht hat. Überhaupt sind die Werke X, 1, S. 167—170 unter
IV. zusammengestellten Leistenotizen den \textls[120]{irregulären} zuzuzählen.

Auf die Siebenteilung der Lemniskate bezieht sich die Leistenotiz Werke
X, 1, S. 160, Art. 1, auf die Fünfteilung (vergl. \textit{Tagebuch} 62. vom 21. März
1797) beziehen sich Leistenotizen, die Werke X, 1, S. 160, Art. 2 und S. 161,
162, Artikel 4, 5 abgedruckt sind.
